%%%%%%%%%%%%%%%%%%%%%%%%%%%%%%%%%%%%%%%%%%%%%%%%%%%%%%%%%%%%%%%%%%%%%%%%%%%
%%%                         JUSTIFICATIVA                                %%%
%%%%%%%%%%%%%%%%%%%%%%%%%%%%%%%%%%%%%%%%%%%%%%%%%%%%%%%%%%%%%%%%%%%%%%%%%%%

\chapter{Justificativa}
\label{ch:justificativa}

Diante do cenário exposto na introdução deste documento, evidencia-se a
necessidade de estratégias pedagógicas que fortaleçam o ensino de Matemática,
especialmente no que se refere ao desempenho dos estudantes em Matemática no
ENEM. Nesse contexto, o Pensamento Computacional, alinhado às diretrizes da
Base Nacional Comum Curricular, destaca-se como uma abordagem relevante
para o desenvolvimento de habilidades para resolução de problemas dos alunos.

Neste sentido, a proposta busca integrar esses referenciais em um ambiente de
aprendizagem real. Assim, espera-se contribuir para a melhoria do processo de
aprendizagem, investigando como a articulação entre Pensamento Computacional e gamificação pode
impactar positivamente o desempenho dos estudantes no ENEM.
